% --- Άσκηση 35410 (35410.tex) ---
\Askhsh{35410}{4}
\begin{ylh}{2}
    \item 4.2. Ανισώσεις 2ου Βαθμού
    \item 6.2. Γραφική Παράσταση Συνάρτησης
\end{ylh}

Δίνεται το τριώνυμο: $T(x) = x^2 - 2a x + a$, με παράμετρο $a \in \mathbb{R}$.

\begin{alist}
\item Να αποδείξετε ότι η διακρίνουσα του τριωνύμου είναι $\Delta = 4a(a - 1)$.
\monades{05}
\item Θεωρούμε την συνάρτηση $f$, που είναι ορισμένη στο $\mathbb{R}$ με τύπο $f(x) = x^2 - 2a x + a$.

Στο καθένα από τα επόμενα σχήματα δίνεται η γραφική παράσταση της συνάρτησης $f$ για διαφορετικές τιμές της παραμέτρου $a$.

\begin{rlist}
\item Για τα δύο πρώτα σχήματα δίνεται ότι η παράμετρος $a \in \{0, 1\}$. Να βρείτε σε ποια τιμή του $a$ αντιστοιχεί το καθένα από τα σχήματα αυτά, δικαιολογώντας την απάντηση σας.
\monades{10}
\item Για το σχήμα 3 να βρείτε τις δυνατές τιμές που μπορεί να πάρει η παράμετρος $a \in \mathbb{R}$, δικαιολογώντας την απάντηση σας.
\monades{10}
\end{rlist}
\end{alist}

\begin{center}
\begin{tikzpicture}[scale=0.8]
\begin{axis}[
    axis lines = middle, xlabel={$x$}, ylabel={$y$},
    xmin=-3, xmax=4, ymin=-1, ymax=5,
    width=7cm, height=6cm,
    ticks=none,
    title={\small Σχήμα 1},
]
